\documentclass[./examen.tex]{subfiles}
\usepackage[first=1, last=20]{lcg}
\newcommand{\random}{\rand\arabic{rand}}
\usepackage{verbatim}
\newcounter{z}
\newcommand\y[2]{
  \loop \ifnum\value{z} < #1
    #2%
    \stepcounter{z}%
  \repeat
}

\begin{document}
\begin{comment}
    
    \question Responder las siguientes problemas:
    
       \end{comment}


%\info{Evaluación Electrotécnia II}{Campo magnético, Inducción magnética, permeabilidad magnética,imanes, líneas de campo magnético, flujo magnético, unidades magnéticas}
\y{11}{
\begin{center}

\makebox[\textwidth]{Nombre y Apellido:\enspace\hrulefill}
\end{center}
\begin{questions}
   
    \question Resolver los siguientes problemas:
     \begin{parts}
     \part[4] Calcular la fuerza obtenida $F_1$ en el émbolo mayor de una prensa hidráulica si en el menor se hacen $F_2=\random N$ y los émbolos circulares tienen triple radio uno del otro $r_1=3 r_2$.
    \part[2]  El área de un pistón en una bomba de fuerza es de $A_1=10cm^2$.¿Qué fuerza $F_2$ se requiere para elevar agua con el pistón hasta una altura de $h=\random m$?
    \part[2]  ¿Qué profundidad de agua sería necesaria para alcanzar la presión absoluta $p=\random kg/cm^2$?s
    \part[2] ¿A qué presión se encuentra un barco Romano del siglo IV en el mar Mediterraneo si está a la profundidad de $h=98m$?
    
    \end{parts}
\end{questions}

}
\end{document}
